\scanpage{30}
\mspage{28}
également les propriétés de rigidité.\inserted{]}

Pour en revenir à l'alinéa précédent, concernant un schéma abélien ``ordinaire'' en car. $p>0$, s'il n'est pas exclu que le $H^1_{\operatorname{DR}}$ admette une stratification canonique fonctorielle, il me semble eependant exclu, malheureusement, que celle-ci provienne d'un cristal \deleted{sur $S$} de modules sur $S$, i.e. que ceci reste vrai en remplaçant $S$ par toute extension infinitésimale, pas néc de car.$p$, \inserted{(tout-au moins} en admettant que la connexion associée à la stratification en question soit la connexion de Gauss-Manin\inserted{)}. En effet, appliquant une telle hypothèse au schéma modulaire $S$ sur $\underline{Z}_p$ précédent, dont on enlèverait les points de Hasse nul d'abord\inserted{, d'où $S'$}, on conclurait sauf erreur que \deleted{sur $S$ tout entier,} la stratification qu'on a en caractéristique nulle sur $H^1_{\operatorname{DR}}$ se prolongerait à $S$ tout entier, car elle se \inserted{``}recollerait'' en un sens évident avec la \deleted{connexion} stratification qu'on aurait \deleted{sur le} \inserted{au dessus du} complété p-adique de $S'$ (ce qui doit impliquer le prolongement sur $S$ de façon assez formelle). Mais alors on aurait en car.$p$ une stratification du $H^1_{\operatorname{DR}}$ au dessus du schéma modulaire tout entier, Hasse nul inclus, ce qui est absurde comme on l'a déjà remarqué. Il faut en conclure, hélas, que si les isocristaux, et le cas échéant les modules stratifiés même en caractéristique \deleted{nulle} $p>0$, ont des chances d'être des outils convenables pour l'étude de familles de variétés abéliennes, celle de cristal elle-même semble irrémédiablement trop fine, même en se restreignant à des familles de VA ``ordinaires'' en car $p>0$. Elle peut tout au mieux \unsure{de|se} prêter au cas d'un schéma de base réduit à un point, spectre d'un corps pas nécessairement parfait, et on peut alors espérer les résultats les plus satisfaisants en se bornant aux VA ordinaires? - Pour que la notion de cristal elle-même puisse être utilisée pour des schémas abéliens sur des bases plus générales, il semble donc qu'il faille imposer aux familles envisagées des restrictions très sérieuses, consistant à imposer la variation infinitésimale du $G(A)$. Cela semble assez proche du point de vue de Serre, qui étudie les variations de VA (éventuellement à multiplication complexe donnée) en imposant à priori l'espace tangent (et l'action de la MC dessus) ...
\margin{29}
